% --- 1. INFORMAZIONI GENERALI ---
\section{Informazioni Generali}
\begin{itemize}
	\item \textbf{Luogo/Piattaforma:} Server Discord
	\item \textbf{Data:} 2026-08-11
	\item \textbf{Orario di inizio:} 14:30
	\item \textbf{Orario di fine:} 15:30
	\item \textbf{Responsabile:} Anton Ricardo Rupi
	\item \textbf{Scriba:} Sofia De Blasi
\end{itemize}

\vspace{0.5cm}
\textbf{Partecipanti presenti:}
\begin{table}[ht]
	\centering
	\renewcommand{\arraystretch}{1.2}
	\begin{tabularx}{\textwidth}{X l}
		\toprule
		\textbf{Nome e Cognome} & \textbf{Matricola} \\
		\midrule
		Sofia De Blasi & 2111014 \\
		Roberto Mariano Doroftei & 2111031 \\
		Filippo Idiometri & 2035617 \\
		Francesco Marcon & 2110990 \\
		Armando Moda Scarati & 2082864 \\
		Anton Ricardo Rupi & 2054304 \\
		\bottomrule
	\end{tabularx}
\\ \end{table}

\vspace{0.5cm}
\textbf{Partecipanti assenti:}
\begin{table}[ht]
	\centering
	\renewcommand{\arraystretch}{1.2}
	\begin{tabularx}{\textwidth}{X l}
		\toprule
		\textbf{Nome e Cognome} & \textbf{Matricola} \\
		\midrule
		Nessuno & - \\
		\bottomrule
	\end{tabularx}
\end{table}

% --- 2. ORDINE DEL GIORNO ---
\section{Ordine del Giorno}
Gli argomenti previsti per la riunione odierna sono:
\begin{enumerate}
	\item Sprint review e retrospective.
	\item Completato lo scheletro delle classi progettate per il backend e frontend.
	\item Conclusa la prima stesura della specifica tecnica.
	\item Dubbi riguardanti il deploy.
	\item Sistemare riferimenti PdP e PdQ con ultimo accesso.
	\item Aggiornare il mvp.
	\item Prima stesura del manuale utente.
	\item Assegnazione dei nuovi ruoli e pianificazione delle attività.
\end{enumerate}

% --- 3. RESOCONTO E DISCUSSIONE ---
\section{Resoconto della Riunione}

\subsection{Sprint review e retrospective}
Durante lo sprint è stato completato lo scheletro delle classi progettate per il backend. 
Analogamente, per il frontend. 
È stata inoltre conclusa la prima stesura della specifica tecnica.
Durante la retrospective è emersa una difficoltà legata alla disponibilità simultanea dei membri del gruppo, 
che ha reso più complesso coordinare alcune attività. La problematica è stata affrontata durante le riunioni, 
cercando di migliorare la comunicazione all'interno del gruppo e aumentando gli aggiornamenti reciproci sull'avanzamento delle attività durante lo sprint.
Un'ulteriore problematica è stata la ridotta responsività della proponente durante il periodo estivo, 
che ha comportato possibili tempi di attesa più lunghi per ottenere risposte a domande o richieste di chiarimento. 
Per evitare che tali attese risultassero bloccanti, 
si è cercato di anticipare il più possibile le domande e le richieste, 
organizzando le attività in modo da poter proseguire autonomamente durante i periodi di attesa.


\subsection{Dubbi riguardanti il deploy}
Chiedere alla proponente se, per il deploy, 
sia necessario predisporre un sito accessibile online oppure se 
sia sufficiente fornire il progetto in formato .zip, 
eseguibile tramite Docker.

\subsection{Sistemare riferimenti PdP e PdQ con ultimo accesso}
Sistemare i riferimenti per PdP e PdQ aggiungendo l'ultimo accesso.

\subsection{Aggiornare il mvp}
Per il completamento dell’MVP rimangono ancora alcune attività da portare a termine.
In particolare, è necessario completare lo sviluppo dei componenti effettivi in Vue.js 
nell’ambito della feature frontend, realizzare i test di integrazione relativi all’infrastruttura 
e completare i test backend con Pytest. Infine, dovranno essere effettuate le 
misurazioni necessarie per il PdQ, considerate parte delle attività infrastrutturali previste per l’MVP.

\subsection{Prima stesura del manuale utente}
Inizia la prima stesura del manuale utente.


La distribuzione del lavoro per l'implementazione è la seguente:
\begin{itemize}
	\item \textbf{Programmatori (due flussi):} Sofia de Blasi si dedicherà allo sviluppo dei componenti effettivi in Vue.js 
	nell’ambito della feature frontend, realizzare i test di integrazione relativi all’infrastruttura e completare i test backend con Pytest. 
	Roberto Mariano Doroftei si dedicherà misurazioni necessarie per il PdQ, considerate parte delle attività infrastrutturali previste per l’mvp.
	\item \textbf{Responsabile:} Anton Ricardo Rupi redigerà il seguente verbale e si occuperà del Diario di Bordo.
	\item \textbf{Amministratore:} Armando Moda Scarati aggiornerà \textbf{PdP} e \textbf{Pdq} aggiungendo l'ultimo accesso per i riferimenti.
	\item \textbf{Progettista:} Francesco Marcon inizierà la stesura del manuale utente.
	\item \textbf{Verificatore:} Filippo Idiometri.
\end{itemize}

% --- 4. DECISIONI PRESE ---
\section{Decisioni Prese}

\renewcommand{\arraystretch}{1.3}
\vspace{-0.3cm}
\noindent
\begin{longtable}{c p{12cm}}
	\toprule
	\textbf{Codice} & \textbf{Decisione} \\
	\midrule
	\endfirsthead

	\toprule
	\textbf{Codice} & \textbf{Decisione} \\
	\midrule
	\endhead

	\midrule
	\multicolumn{2}{r}{\textit{Continua nella pagina successiva}} \\
	\midrule
	\endfoot

	\bottomrule
	\endlastfoot

	\textbf{VI-26.1} & {Sistemare i riferimenti per PdP e PdQ aggiungendo l'ultimo accesso.} \\
	\textbf{VI-26.2} & {Aggiornare il mvp.} \\
	\textbf{VI-26.3} & {Prima stesura del manuale utente.} \\
	
\end{longtable}

% --- 5. AZIONI (TODO) ---
\section{Azioni da Svolgere (To-Do)}

\renewcommand{\arraystretch}{1.3}
\begin{longtable}{p{2.3cm} c p{4.5cm} p{3cm} l}
	\toprule
	\textbf{Codice} & \textbf{Rif. Decisione} & \textbf{Descrizione Task} & \textbf{Assegnatario} & \textbf{Scadenza} \\
	\midrule
	\endfirsthead

	\toprule
	\textbf{Codice} & \textbf{Rif. Decisione} & \textbf{Descrizione Task} & \textbf{Assegnatario} & \textbf{Scadenza} \\
	\midrule
	\endhead

	\midrule
	\multicolumn{5}{r}{\textit{Continua nella pagina successiva}} \\
	\midrule
	\endfoot

	\bottomrule
	\endlastfoot

	\ghissue{359} & - & Stesura del verbale interno odierno (VI N. 26) & Anton Ricardo Rupi & 2026-08-18 \\[4pt]
	\ghissue{360} & - & Stesura del Diario di Bordo & Anton Ricardo Rupi & 2026-08-14 \\[4pt]
	\ghissue{361} & VI-26.1 & Aggiornamento PdP & Armando M. Scarati & 2026-08-18 \\[4pt]
	\ghissue{362} & VI-26.1 & Aggiornamento PdQ & Armando M. Scarati & 2026-08-18 \\[4pt]
	\ghissue{363} & - & Invio email a Zucchetti & Francesco Marcon & 2026-08-18 \\[4pt]
	\ghissue{364} & VI-26.2 & Sviluppo componenti vue.js effettivi & Sofia de Blasi & 2026-08-18 \\[4pt]
	\ghissue{365} & VI-26.2 & Aggiunta test infra test integrazione & Sofia de Blasi & 2026-08-18 \\[4pt]
	\ghissue{366} & VI-26.2 & Aggiunta test backend test pytest & Sofia de Blasi & 2026-08-18 \\[4pt]
	\ghissue{367} & VI-26.3 & Prima stesura manuale utente & Francesco Marcon & 2026-08-18 \\[4pt]
	\ghissue{368} & VI-26.2 & Effettuare le misurazioni necessarie per il PdQ & Roberto M. Doroftei & 2026-08-18 \\[4pt]
	
\end{longtable}

% --- 6. PROSSIMA RIUNIONE ---
\section{Prossima Riunione}
\begin{itemize}
	\item \textbf{Data:} 2026-08-18
	\item \textbf{Ora:} 14:30
	\item \textbf{OdG preliminare:}
	\begin{itemize}
		\item Stesura verbale esterno.
		\item Revisione funzionamento mvp.
		\item Revisione manuale utente.
		\item Pianificazione del prossimo sprint e assegnazione dei ruoli.
	\end{itemize}
\end{itemize}